\documentclass{article}
\usepackage{geometry}
\usepackage{graphicx} % Required for inserting images
\usepackage{amsmath, amsthm, amssymb}
\usepackage{parskip}
\newgeometry{vmargin={15mm}, hmargin={24mm,34mm}}
\theoremstyle{definition} 
\newtheorem{definition}{Definition}

\newtheorem{theorem}{Theorem}[section]
\newtheorem{lemma}[theorem]{Lemma}
\newtheorem{proposition}[theorem]{Proposition}
\newtheorem{corollary}{Corollary}[theorem]


\title{MATH210B Homework 8}
\date{March 2025}
\author{Boran Erol}

\begin{document}

\maketitle

\section{Problem 1}

Let $ \zeta $ be a primitive p-th root of unity.

Recall that $ [\mathbb{Q}(\zeta): \mathbb{Q} ] = p - 1 $.

Notice that $ \mathbb{Q}(\zeta + \zeta^{-1}) \subseteq \mathbb{R} $ and recall that 
$ [\mathbb{Q}(\zeta): \mathbb{Q}(\zeta + \zeta^{-1}, \zeta \zeta^{-1}) ] \leq 2 $.

Since $ \mathbb{Q}(\zeta) $ is not contained in $ \mathbb{R} $, then this is exactly $ 2 $.

By the tower law, $ [\mathbb{Q}(\zeta + \zeta^{-1}) : \mathbb{Q} ] = \frac{p - 1}{2} $.

(Notice $ \zeta \zeta^{-1} = 1 $ so it doesn't contribute anything.)


\section{Problem 2}

Notice that $ \mathbb{Q}(\sqrt{2} + \sqrt{3}) $ contains $ \sqrt{6} $ as well.
Thus, $ \mathbb{Q}(\sqrt{2} + \sqrt{3}) =  \mathbb{Q}(\sqrt{2} + \sqrt{3}, \sqrt{2}\sqrt{3}) $.

Since the minimal polynomial of $ \sqrt{2} + \sqrt{3} $ has degree 4, it's a degree 4 extension.

Notice that $ \mathbb{Q}(\sqrt{2} + \sqrt{3}) $ is also degree 4. This is because $ \sqrt{3} $ is not
in $ \mathbb{Q}(\sqrt{2}) $, which is proven at the end. Thus, these two extensions are equal.
However, $ \mathbb{Q}(\sqrt{2} + \sqrt{3}) $ is separable since it is the splitting field of $ (x^{2} - 2)(x^{2} - 3) $,
so we're done.

\begin{lemma}
    $ \sqrt{3} \notin \mathbb{Q}(\sqrt{2}) $.
\end{lemma}
\begin{proof}
    Assume it is. Then, $ \sqrt{3} = a + b \sqrt{2} $ for some $ a,b \in Q $.
    Squaring both sides gives $ 3 = a^{2} + 2ab \sqrt{2} + 2b^{2} $.

    Then, $ ab = 0 $. Assume first that $ a = 0 $. Then, $ 3 = 2b^{2} $.
    Assume $ b = \frac{p}{q} $ where $ p,q $ are coprime. Then, $ 3q^{2} = 2p^{2} $,
    so $ 2 $ divides $ q^{2} $, so $ 2 $ divides $ q $, so $ 4 $ divides $ q^{2} $,
    so $ 2 $ divides $ p $, which is a contradiction.

    Assume now $ b = 0 $. Then, $ a = \sqrt{3} $ which is not true since $ 3 $ isn't rational.
\end{proof}

\section{Problem 4}

Recall that every finite extension of a finite field is normal, since it's a splitting field.

Then, we need to find the smallest $ d $ such that $ 7 $ divides $ 5^{d} - 1 $.

This is equivalent to finding the multiplicative order of $ 5 $ in $ (\mathbb{Z} / 7 \mathbb{Z})^{\times} $,
which is $ 6 $. Thus, $ d = 6 $.

\section{Problem 5}

Let $ char(F) = p > 0 $.

Let $ E/F $ be an extension such that $ [E : F] $ and $ p $ are coprime.

Let $ E = F(\alpha_{1}, \ldots, \alpha_{n}) $.

Notice that the degree of $ m_{\alpha_{i}} $ is coprime to $ p $ since it divides $ [E:F] $.
Let this degree be $ n_{i} $.

Then, the derivative of $ m_{\alpha_{i}} $ is non-zero in $ F $ since the leading coefficient is monic
and so the leading coefficient of the derivative is $ n_{i} \mod p $ which is nonzero. Thus, $ E $
is generated by separable elements over $ F $ and thus is separable over $ F $.

\section{Problem 6}

Let $ E/F $ be a finite field extension.

Let $ \alpha, \beta \in E $ be separable over $ F $.

Then, $ F(\alpha, \beta) $ is separable over $ F $ since it's generated by separable elements
using Corollary 3.5.20 in the notes. But then, $ \alpha + \beta $ and $ \alpha \beta $ is in $ E $
since this is a subfield of $ E $. Notice that this also contains $ \alpha^{-1} $.

Let's now find the maximal separable extension of $ F(X) / F(X^{4} + X^{2}) $, where $ F $ is
a field of characteristic $ 2 $.

Let's first show that $ f(t) = t^{4} + t^{2} + (X^{4} + X^{2}) $ is the minimal polynomial of $ X $.
Clearly, $ X $ is a root, so let's show that it's irreducible.

Notice that 

$ f(t) = t^{4} + t^{2} + (X^{4} + X^{2}) = (t + X)^{4} + (t + X)^{2} = (t + X)^{4} + (t + X)^{2} = (t + X)^{2}((t + X)^{2} + 1) $.
But the roots of this are $ X $ and $ X + 1 $, so the polynomial doesn't split over $ F(X^{4} + X^{2}) $.
Thus, $ F(X) $ is an extension of degree $ 4 $. Notice that $ f $ isn't separable since it's
derivative is zero.

Also notice that $ X^{2} $ is separable since it's the root
of $ t(t + 1) + (X^{4} + X^{2}) $, which is separable and irreducible using the same argument as above.
Thus, the maximal separable subextension is $ F(X^{2}) $.

\section{Problem 7}

Let $ E = F(\alpha_{1}, \ldots, \alpha_{n}) $.

By Homework 9 Problem 1 (where I didn't use this problem), we have that
the minimal polynomial of every element contains exactly one root. But then the extension should
map each $ \alpha_{i} $ to itself, so it acts like the identity on $ E $, and therefore there's
exactly one extension.

Now, let's do the backwards direction. Assume there is at most one extension. Then, every
minimal polynomial should be of the form $ (x - \alpha_{1})^{m}  $ for some $ m $ since otherwise
we would have two extension by mapping one root to another root as in Proposition 3.2.4.

We thus conclude the proof.

\section{Problem 8}

Recall that in Problem 1 we showed that $ \mathbb{Q}[\zeta_{p} + \zeta_{p}^{-1}] $ was a degree
$ \frac{p - 1}{2} $. This extension is a subfield of $ \mathbb{Q}[\zeta] $. Since this
is an Abelian extension, we know that subfields are also Galois extensions. Moreover, we know that
$ n $ and $ m $th roots of unity, for coprime $ m,n $ are linearly disjoint over $ \mathbb{Q} $.
Then, adding $ \zeta_{7} + \zeta_{7}^{-1} $ and $ \zeta_{11} + \zeta_{11}^{-1} $ gives us the desired result.

\section{Problem 9}


Let $ f(x) = x^{6} + 3 $.

Let $ \alpha $ be a root of $ f $.

Then, $ \alpha^{3} = \pm \sqrt{3}i $.

Notice that $ \frac{1 + \pm \alpha^{3}}{2} $ is a primitive 6th root of unity.

Thus, the splitting field is generated by $ \alpha $.

Since the action of the Galois group is transitive, $ \alpha $ can go to any other root of $ x^{6} - 3 $,
but then the Galois group is cyclic of order $ 6 $.

\section{Problem 10}





\end{document}
